\documentclass[conference]{IEEEtran}
\usepackage{cite}
\usepackage{amsmath,amssymb,amsfonts}
\usepackage{booktabs}
\usepackage{array}
\usepackage{multirow}
\usepackage{xcolor}
\usepackage{url}
\usepackage{microtype}
\usepackage{balance}
\setlength{\tabcolsep}{3pt}
\newcolumntype{L}[1]{>{\raggedright\arraybackslash}p{#1}}
\newcommand{\evotrace}{\textsc{Evotrace}}
\newcommand{\designonly}{\textsc{Design-Only}}

\begin{document}

\title{How Organisms Evolve: \\
\evotrace{} for Multiscale Causal Evidence from Variation to Complex Traits}

\author{\IEEEauthorblockN{Four-Stage Research Workflow}
\IEEEauthorblockA{\textit{Survey - Idea - Experiment Design - Author}\\
\textnormal{\designonly{} research artifact}\\
Evolutionary biology research design}}

\maketitle

\begin{abstract}
Evolution is inherited change in populations across generations, not a property of an individual organism alone. Natural selection is central to adaptive evolution, but observed evolutionary change can also reflect mutation, recombination, genetic drift, gene flow, demographic history, and environmental heterogeneity. Complex organs introduce an additional demand: a population-level trajectory must be connected to genotype, developmental regulation, phenotype, and fitness without treating correlation as causal proof. This design-only paper introduces EVOTRACE, Evolutionary Trait Reconstruction and Causal Evidence, a five-stage framework for this problem. EVOTRACE fixes a lineage and environmental passport, validates heritable variation, compares population-process models, tests genotype-to-development-to-trait paths, and evaluates predictions on held-out lineages or interventions. The proposal joins time-series genomics, ecology, common-environment assays, developmental atlases, and controlled perturbations while retaining explicit alternative explanations. It formalizes uncertainty from drift, migration, ancestry, stage alignment, and historical reconstruction; it also defines promotion rules for descriptive, selection-supported, mechanistically supported, and historical claims. The conclusion is that a persuasive account of how organisms evolve is a linked causal evidence chain rather than an endpoint narrative. The paper advances no new evolutionary data, selection coefficients, or organ-evolution result.
\end{abstract}

\begin{IEEEkeywords}
evolution, natural selection, genetic drift, gene flow, experimental evolution, evolutionary developmental biology, gene regulatory networks, causal inference
\end{IEEEkeywords}

\section{Introduction}

Evolutionary change occurs when inherited variants change in frequency across generations. This concise statement conceals several interacting processes. New variants arise through mutation, recombination, duplication, rearrangement, and, in some lineages, gene transfer. Their frequencies then reflect differential survival and reproduction, random sampling, migration, demographic history, and the mapping from genotype and environment to phenotype. Natural selection is the principal mechanism by which traits that increase relative reproductive success can spread, but a frequency shift alone does not identify selection.

The need to separate mechanisms is especially clear in genomic time series. A rise in an allele can result from positive selection, a population bottleneck, linkage to another selected locus, migration, or a sampling artifact. Likewise, a difference in gene expression between species can reflect a developmental-stage mismatch, cell-composition difference, neutral regulatory turnover, or a causal developmental change. A general account of evolution must expose these alternatives rather than allow an adaptive story to become the default interpretation.

Experimental evolution supplies unusually strong leverage. Wiser, Ribeck, and Lenski measured fitness trajectories of 12 \emph{Escherichia coli} populations over 50,000 generations and found continued gains consistent with a power-law model involving diminishing-returns epistasis and clonal interference~\cite{wiser2013}. The result demonstrates that adaptation can be measured as a time-resolved process under controlled conditions; it does not imply that every lineage, environment, or trait follows the same trajectory. Blount \emph{et al.} further showed that a key innovation in an experimental bacterial population required historical potentiation, illustrating why the observation of a beneficial phenotype alone cannot establish inevitability~\cite{blount2008}.

Complex traits require a second scale of explanation. The developmental genetics of eyes, limbs, and other structures indicates that novelties often emerge through modification and redeployment of pre-existing regulatory circuits~\cite{shubin2009}. Hierarchical gene regulatory networks (GRNs) provide one account of why changes at different regulatory levels can have distinct phenotypic consequences~\cite{erwindavidson2009}. This evidence motivates a framework that connects population genetics to development without assuming that a shared gene automatically explains a shared structure.

This paper introduces \evotrace{}, Evolutionary Trait Reconstruction and Causal Evidence. The framework is designed for a focal lineage, trait, and environment, with optional phylogenetic and experimental replication. Its core principle is simple: a claim about ``how a trait evolved'' is promoted only after the relevant evidence chain is measured. The report is \designonly{}; it presents no new genomic observations, interventions, historical reconstructions, or organismal fitness data.

\section{Mechanisms of Evolution and Their Evidentiary Boundaries}

\subsection{Variation, transmission, and sorting}

The raw material of evolution is heritable variation, but the origin of a variant and the process that changes its frequency should not be conflated. Mutation produces new alleles; recombination reshuffles existing variants; duplication and transposable elements can change genomic architecture; and standing variation can permit rapid response when environments change~\cite{barrettschluter2008,kaessmann2010}. Selection changes expected transmission through relative fitness. Drift produces stochastic changes that are particularly important when effective population size is small. Gene flow introduces alleles and can either oppose local adaptation or supply useful variants. These processes can operate simultaneously.

\begin{table*}[t]
\caption{Evolutionary processes, observable signals, and the alternative explanations that must remain in the analysis.}
\label{tab:forces}
\centering
\footnotesize
\begin{tabular}{L{0.15\textwidth} L{0.23\textwidth} L{0.28\textwidth} L{0.24\textwidth}}
\toprule
\textbf{Process} & \textbf{Population-level prediction} & \textbf{Minimum evidence} & \textbf{Alternative to test}\\
\midrule
Mutation/recombination & New alleles or haplotypes enter the population & pedigree or time-series sequence data, error model, and recombination uncertainty & sequencing artifact, recurrent mutation, or unobserved ancestry\\
Natural selection & Frequency or trait change covaries with relative fitness in a declared environment & replicated fitness/performance data, ecological covariates, and demographic null model & drift, linked selection, population structure, or migration\\
Genetic drift & Replicate trajectories vary stochastically around a neutral expectation & effective population size, neutral references, temporal samples & unmodelled selection or sampling bias\\
Gene flow & Ancestry or allele sharing follows migration/connectivity pathways & geography, ancestry/admixture, migration model, and sampling frame & shared ancestral polymorphism or convergent mutation\\
Developmental regulation & Candidate variants alter cell state, expression, or developmental timing linked to a trait & matched stages, cell states, spatial expression, and perturbation/orthogonal evidence & stage mismatch, cell-composition shift, or reactive expression\\
Ecological feedback & Trait effects differ across environments and alter subsequent selection & common-environment or reciprocal exposure plus fitness endpoints & laboratory-only performance or unmeasured environmental confounding\\
\bottomrule
\end{tabular}
\end{table*}

Table~\ref{tab:forces} maps mechanisms to evidence, not to narrative labels. Molecular signatures can be informative when matched to a well-specified demographic model~\cite{nielsen2005}, but no single summary statistic is a universal test of adaptation. The neutral theory remains a crucial reference because it establishes that random drift can create molecular change without adaptive cause~\cite{kimura1983}. Conversely, neutral models do not make selection untestable; they define the competing prediction that data must exceed.

\subsection{Complex traits and evolutionary novelty}

The evolution of a complex organ should be decomposed into intermediate questions: Which precursor structures or cell types existed? Which modules were modified, co-opted, duplicated, or lost? How did regulatory deployment change in space and time? What performance and fitness consequences followed in a particular ecology? The resulting account may include gradual phenotypic change, network-level rewiring, or multiple coordinated changes, but it should not begin by assuming that the present organ was the target of a directed process.

At larger evolutionary scales, the origin and sorting of variation need to be studied together rather than treating diversity, morphology, and function as interchangeable outcomes~\cite{jablonski2017}. A replicated natural example can connect ecology and developmental regulation: Favé \emph{et al.} examined responses of ant populations across climate-isolated Sky Islands, illustrating how developmental-network and phenotypic evidence can be assessed alongside population history~\cite{fave2015}. Such case studies motivate the multi-layer design below but do not substitute for replication in a new focal trait.

Deep homology provides a useful testable concept. Shubin, Tabin, and Carroll describe cases in which early metazoan regulatory mechanisms underlie independently evolved structures, including eyes and limbs~\cite{shubin2009}. The claim predicts that comparison of homologous cell programs and regulatory connections can reveal conserved generative components, even where mature anatomy differs. It does not erase independent innovation: a conserved toolkit may be deployed in different tissues, at different times, and under different selection regimes.

Placental and pregnancy-related traits make the same point from another direction. Vertebrate reproductive modes and placental structures show repeated origin, diversification, and potential conflict between maternal and fetal interests~\cite{crespsemen2004}. A research design must distinguish repeated ecological association, shared ancestry, gene recruitment, regulatory co-option, and experimentally measured developmental function. A claim that a particular regulatory change ``created'' a placenta is much stronger than the observation that it is expressed in a placenta.

\begin{table}[t]
\caption{Claim ladder for an evolutionary explanation.}
\label{tab:claims}
\centering
\footnotesize
\begin{tabular}{L{0.28\columnwidth} L{0.29\columnwidth} L{0.29\columnwidth}}
\toprule
\textbf{Claim level} & \textbf{What the claim says} & \textbf{Evidence threshold}\\
\midrule
Descriptive trajectory & Allele, haplotype, trait, or expression measure changed in a defined population/time frame & quality-controlled repeated observations and uncertainty interval\\
Selection-supported & A change is better predicted by an environment-dependent selection model than declared alternatives & demographic, migration, and drift alternatives plus replication or held-out prediction\\
Mechanistically supported & A genomic/regulatory candidate affects a measured intermediate and trait in the predicted direction & directional perturbation or strong orthogonal evidence, stage-matched controls, and off-target checks\\
Historical explanation & A proposed path from ancestral state to present trait is the best supported among explicit alternatives & phylogenetic/ancestral uncertainty, comparative replication, and a mechanism-compatible route\\
\bottomrule
\end{tabular}
\end{table}

Table~\ref{tab:claims} prevents an association from being upgraded by language alone. A study may legitimately report a descriptive trajectory while remaining agnostic about the selection mechanism; it may report a molecular mechanism without claiming that the mechanism was the historical cause of a lineage divergence. This separation makes negative and inconclusive results scientifically useful.

\section{The \evotrace{} Framework}

\subsection{Evolutionary passport}

EVOTRACE starts by freezing a target claim. Let the evolutionary estimand be
\begin{equation}
\mathcal{Q}=\left(L,Z,E,T,G,P,\mathcal{M},\mathcal{C}\right),
\label{eq:passport}
\end{equation}
where $L$ is the lineage/population boundary, $Z$ the trait and life-stage definition, $E$ the environmental frame, $T$ the generation or calendar-time window, $G$ the genomic unit, $P$ the sampling and measurement protocol, $\mathcal{M}$ the competing population-process models, and $\mathcal{C}$ the intended claim level. Equation~\eqref{eq:passport} makes the scope of a conclusion inspectable. A result cannot be interpreted as a general account of adaptation if the trait, environment, lineage, or comparator is unspecified.

The passport also records which data are confirmatory and which are exploratory; the reference genome and annotation release; ancestry and relatedness rules; environmental exposure; developmental-stage alignment; phenotype definition; fitness endpoint; and exclusion criteria. It is locked before inspecting primary evolutionary model fits. This avoids selecting a convenient population, trait measure, or neutral reference after observing a result.

\begin{table}[t]
\caption{EVOTRACE stages and their fail-fast decisions.}
\label{tab:stages}
\centering
\footnotesize
\begin{tabular}{L{0.13\columnwidth} L{0.36\columnwidth} L{0.42\columnwidth}}
\toprule
\textbf{Stage} & \textbf{Required output} & \textbf{Fail-fast decision}\\
\midrule
E0 Scope & Passport, comparator, causal target, and claim level & Stop if lineage, environment, trait, or time scale is undefined.\\
E1 Variation & QC-ed genotype, ancestry, phenotype, and developmental ledger & Stop if measurement error, relatedness, or stage alignment cannot be assessed.\\
E2 Transmission & Competing selection, drift, mutation, and migration predictions & Stop promotion if a demographic/neutral alternative is absent.\\
E3 Mechanism & Genotype-to-intermediate-to-trait evidence map & Stop at association if directionality, controls, or plausible intermediate are missing.\\
E4 Evaluation & Held-out calibration, sensitivity, and independent replication decision & Downgrade inference if the account does not generalize.\\
\bottomrule
\end{tabular}
\end{table}

Table~\ref{tab:stages} describes the five evidence gates. The gates are not a demand that every study perform every assay. They clarify what a study has established and what additional evidence would be needed for a stronger claim. A comparative atlas may complete E1 and motivate E3 experiments; a time-series evolution experiment may strongly support E2 without resolving the developmental basis of a trait.

\begin{figure*}[t]
\centering
\fbox{\begin{minipage}{0.94\textwidth}
\centering
\vspace{0.45em}
\begin{tabular}{@{}c@{\hspace{0.75em}}c@{\hspace{0.75em}}c@{\hspace{0.75em}}c@{\hspace{0.75em}}c@{}}
\fbox{\parbox[c][2.4em][c]{0.15\textwidth}{\centering\footnotesize\textbf{E0 Scope}\\Lineage, trait, ecology}} &
$\Longrightarrow$ &
\fbox{\parbox[c][2.4em][c]{0.15\textwidth}{\centering\footnotesize\textbf{E1 Variation}\\Inherited variants and phenotypes}} &
$\Longrightarrow$ &
\fbox{\parbox[c][2.4em][c]{0.15\textwidth}{\centering\footnotesize\textbf{E2 Transmission}\\Selection, drift, and flow models}}\\\noalign{\vskip1.15em}
\fbox{\parbox[c][2.4em][c]{0.15\textwidth}{\centering\footnotesize\textbf{E3 Mechanism}\\Developmental intermediates}} &
$\Longrightarrow$ &
\fbox{\parbox[c][2.4em][c]{0.15\textwidth}{\centering\footnotesize\textbf{E4 Evaluation}\\Prediction and replication}} &
$\Longrightarrow$ &
\fbox{\parbox[c][2.4em][c]{0.30\textwidth}{\centering\footnotesize\textbf{Versioned explanation}\\Descriptive, selection-supported, mechanistic, or historical claim}}
\end{tabular}
\vspace{0.65em}
\end{minipage}}
\caption{EVOTRACE links population-level transmission to developmental mechanism and evaluation. Failure at a stage narrows the claim; it does not convert a remaining association into causal evidence.}
\label{fig:evotrace}
\end{figure*}

Fig.~\ref{fig:evotrace} represents the intended information flow. E0 is upstream: neither selection coefficients nor gene expression can be interpreted without a declared target population and environment. E1 and E2 establish whether a population process needs explanation. E3 asks whether a developmental or physiological path mediates a trait, and E4 tests whether the linked account predicts new observations. The framework's final output is an evidence-labeled explanation, not a predetermined adaptive narrative.

\section{Population Transmission: Competing Models Rather Than One Statistic}

\subsection{Fitness, selection, and drift}

For haplotypes $i=1,\ldots,K$ with frequency $p_{i,t}$ and relative fitness $w_{i,t}$, mean fitness at time $t$ is
\begin{equation}
\overline{w}_{t}=\sum_{i=1}^{K}p_{i,t}w_{i,t}.
\label{eq:meanfitness}
\end{equation}
Equation~\eqref{eq:meanfitness} defines the normalization required before comparing transmission. Relative fitness must correspond to a preregistered viability, fertility, mating, or composite life-history endpoint in the declared environment. A laboratory growth rate may be an appropriate endpoint for a microbial experiment but not necessarily for a field vertebrate trait.

One schematic state transition that keeps multiple processes visible is
\begin{equation}
p_{i,t+1}=(1-m_t)\frac{p_{i,t}w_{i,t}}{\overline{w}_{t}}+m_t p^{(M)}_{i,t}+u_{i,t}+\eta_{i,t},
\label{eq:transition}
\end{equation}
where $m_t$ is migration proportion, $p^{(M)}_{i,t}$ the migrant-pool frequency, $u_{i,t}$ net mutation/recombination input under the chosen representation, and $\eta_{i,t}$ a stochastic term. Equation~\eqref{eq:transition} is a design template rather than a universal biological model. It forces selection, migration, mutation, and randomness to be explicit and allows data to support a simplified special case only after testing its assumptions.

Under an idealized diploid Wright--Fisher neutral approximation, drift variance is
\begin{equation}
\mathrm{Var}\left(\Delta p_t\right)=\frac{p_t(1-p_t)}{2N_e},
\label{eq:drift}
\end{equation}
where $N_e$ is effective population size. Equation~\eqref{eq:drift} provides a baseline for expected random fluctuation, not a substitute for estimating demography. Inference should propagate uncertainty in $N_e$, sample time, genotyping, and linkage; otherwise a narrow neutral envelope can artificially make selection appear certain.

For a biallelic trajectory, an environment-aware selection contrast can be written as
\begin{equation}
\mathrm{logit}(p_{t+1})-\mathrm{logit}(p_t)=s_0+\mathbf{x}_t^{\mathsf{T}}\boldsymbol{s}_E+\epsilon_t,
\label{eq:logit}
\end{equation}
where $s_0$ is a baseline directional component, $\mathbf{x}_t$ contains measured environmental exposures, $\boldsymbol{s}_E$ their interaction coefficients, and $\epsilon_t$ represents remaining process and observation uncertainty. Equation~\eqref{eq:logit} is evaluated against neutral, demographic, and migration alternatives using held-out times or populations. An environmental correlation is not taken as proof of selection when population structure can generate the same pattern.

\begin{table*}[t]
\caption{Sampling and measurement architecture for a first EVOTRACE implementation.}
\label{tab:sampling}
\centering
\footnotesize
\begin{tabular}{L{0.17\textwidth} L{0.25\textwidth} L{0.27\textwidth} L{0.22\textwidth}}
\toprule
\textbf{Evidence layer} & \textbf{Design element} & \textbf{Recorded controls} & \textbf{Primary use}\\
\midrule
Population genomics & Replicated populations across an environmental gradient plus temporal samples where possible & depth, genotype likelihood, relatedness, reference version, recombination uncertainty & ancestry, allele/haplotype trajectories, demographic alternatives\\
Ecology and fitness & Common-environment and reciprocal-environment assays; preregistered survival/fertility/performance endpoint & batch, density, age, sex, exposure duration, observer blinding & genotype/trait-by-environment and relative-fitness estimates\\
Phenotype & Repeated trait measurement at matched life stage and condition & calibration, inter-rater reliability, imaging protocol, missingness & trait trajectory and genotype-to-trait link\\
Developmental regulation & Matched embryological/cell-state time course across focal and comparison lineages & stage atlas, tissue purity, cell composition, library and processing batch & candidate intermediate and GRN deployment\\
Functional test & Allele, regulatory, or pathway perturbation with rescue/orthogonal evidence where permitted & off-target assay, sham/control, independent construct, viability monitoring & directional causal mechanism\\
History/comparison & Calibrated phylogeny, outgroups, fossils or archived samples when applicable & topology/clock uncertainty, character coding, ascertainment rule & alternative historical pathways and convergence tests\\
\bottomrule
\end{tabular}
\end{table*}

\subsection{Replicates and historical contingency}

Replication is needed at the level relevant to the claim. Independent culture lines test repeatability under a controlled environment; independent natural populations test response under related ecological conditions; phylogenetically separate lineages test recurrence across history. Experimental evolution is especially valuable because it can manipulate starting states and environments while measuring fitness and genomes over time~\cite{kawecki2012}. Yet controlled replicate lines should not be used to erase the roles of rare mutations, epistasis, and contingency that emerge under repeated evolution.

Blount \emph{et al.} provide a concrete historical-contingency example: a later capability was accessible only after prior potentiating changes~\cite{blount2008}. EVOTRACE therefore separates a mutation's immediate effect from its effect conditional on genotype and environment. Replaying an evolutionary process may reveal predictable phenotypic patterns, different molecular solutions, or both~\cite{courtier2015}. A prediction of convergence is not a prediction of identical nucleotide substitutions.

\section{From Genotype to Complex Trait}

\subsection{Trait, development, and mediation}

A useful genotype-to-trait model retains intermediate molecular and developmental states. Let $A$ denote a candidate allele or regulatory state, $R$ a measured regulatory/developmental intermediate, $Z$ the trait, and $\mathbf{c}$ confounders such as ancestry, environment, and stage. The trait association may be represented as
\begin{equation}
Z=\alpha_Z+\beta_A A+\beta_R R+\boldsymbol{\gamma}^{\mathsf{T}}\mathbf{c}+e_Z.
\label{eq:trait}
\end{equation}
Equation~\eqref{eq:trait} separates a direct association coefficient $\beta_A$ from the intermediate association $\beta_R$, but it is not a causal proof. The covariate set must include ancestry and developmental-stage controls. Without matched stages, a species difference in $R$ could simply reflect comparison of different points on a developmental trajectory.

The candidate intermediate is modeled as
\begin{equation}
R=\alpha_R+\theta_A A+\boldsymbol{\delta}^{\mathsf{T}}\mathbf{c}+e_R.
\label{eq:mediator}
\end{equation}
Equation~\eqref{eq:mediator} defines the $A\rightarrow R$ link that is needed before discussing mediation. A planned causal pathway requires a directionally interpretable intervention, natural experiment, or orthogonal functional evidence. The product $\theta_A\beta_R$ can summarize a mediation-compatible component only after the model assumptions, measurement time order, and unmeasured-confounding risk are stated.

GRNs are useful because they identify possible intermediates and network levels. Erwin and Davidson argue that network hierarchy can make some perturbations broadly pleiotropic and others more modular~\cite{erwindavidson2009}. A cis-regulatory change may alter timing or tissue-specific deployment with fewer collateral effects than a coding change in some contexts, but this is a hypothesis to test rather than a universal rule~\cite{stern2008}. The experiment should therefore measure regulatory activity, cell state, and phenotype in the same developmental context.

\begin{table*}[t]
\caption{Evidence needed to explain a complex trait without an endpoint-only narrative.}
\label{tab:complex}
\centering
\footnotesize
\begin{tabular}{L{0.17\textwidth} L{0.25\textwidth} L{0.28\textwidth} L{0.21\textwidth}}
\toprule
\textbf{Question} & \textbf{Evidence that supports it} & \textbf{Strong alternative explanation} & \textbf{Required EVOTRACE response}\\
\midrule
Did a trait change? & homologous-stage anatomy/phenotype measurements with uncertainty & measurement protocol or stage difference & preregister trait ontology and blind measurement\\
Did change affect fitness? & survival, fecundity, mating, or performance measured in matched environment & proxy performance unrelated to lifetime reproduction & include ecological exposure and fitness endpoint\\
Was a candidate genetic change involved? & time/lineage association plus ancestry-aware genotype-to-trait evidence & linkage, shared ancestry, or reactive expression & fine-map, model ancestry, and test alternatives\\
Did regulation/development mediate the trait? & spatial and temporal intermediate plus directional perturbation/rescue & stage/cell-composition mismatch or pleiotropic toxicity & matched atlas, controls, viability assessment, orthogonal assay\\
Was the historical route repeated or contingent? & explicit ancestral alternatives and phylogenetic/experimental replication & one reconstructed route or endpoint convergence & report uncertainty and test multiple lineages/start states\\
\bottomrule
\end{tabular}
\end{table*}

Table~\ref{tab:complex} shows why a complex trait cannot be ``explained'' by finding a gene expressed in the mature organ. An eye, for example, requires a defined visual phenotype, relevant cell and developmental stage, environmental and fitness context, and evidence that candidate regulatory changes alter a plausible intermediate. A placenta-related trait requires equivalent attention to maternal tissue, fetal tissue, developmental timing, and reproductive performance. The framework thereby treats complexity as a linked causal question, not as an argument from incredulity.

\subsection{Causal contrasts}

When an intervention is appropriate and ethically permitted, EVOTRACE specifies its estimand in potential-outcome form:
\begin{equation}
\tau_Z=\mathbb{E}\left[Z\left(A=1\right)-Z\left(A=0\right)\right],
\label{eq:ate}
\end{equation}
where $\tau_Z$ is the average effect of the candidate state $A$ on trait $Z$ in the declared experimental population and environment. Equation~\eqref{eq:ate} requires an intervention, exchangeability argument, or justified natural experiment; it cannot be estimated from an observational association alone. For regulatory candidates, designs may use allele swaps, enhancer perturbation, dosage series, rescue, and independent constructs where feasible. The aim is not to recreate a whole historical lineage in the laboratory, but to test a predicted causal link in the proposed pathway.

\section{Evaluation, Sensitivity, and Reproducibility}

\subsection{Competing explanation score}

Each candidate population-process model $M_j$ is evaluated by predictive performance on a held-out unit $h$ (population, time point, lineage, or batch). A generic log predictive score is
\begin{equation}
\mathrm{LPS}(M_j)=\sum_{h=1}^{H_{\mathrm{test}}}\log p\left(D_h\mid D_{-h},M_j\right),
\label{eq:lps}
\end{equation}
where $D_h$ is held-out data and $D_{-h}$ is the retained training data. Equation~\eqref{eq:lps} favors models that forecast new observations rather than simply fit a selected historical dataset. Selection-plus-environment, neutral demographic, migration, and mixed models are all entered before ranking. A higher score is conditional evidence for the stated model set, not an absolute proof that no other mechanism exists.

Sensitivity is reported across a predeclared specification set $a$, including effective population size, migration prior, recombination map, trait extraction, developmental alignment, and ancestral-state model. For an estimate $\widehat{\theta}(a)$, define
\begin{equation}
\Delta(a,a')=\frac{\left|\widehat{\theta}(a)-\widehat{\theta}(a')\right|}{\max\left(\left|\widehat{\theta}(a)\right|,\epsilon\right)},
\label{eq:sensitivity}
\end{equation}
where $a'$ is a plausible alternative and $\epsilon>0$ prevents a degenerate denominator. Equation~\eqref{eq:sensitivity} makes model dependence measurable. Large sensitivity is not automatically a failure; it is evidence that a claim should be narrowed or that the next experiment should target the unstable assumption.

For occurrence, phenotype category, or directional outcome predictions with probability $\pi_r$ and held-out outcome $y_r$, calibration is assessed with
\begin{equation}
\mathrm{BS}=\frac{1}{R}\sum_{r=1}^{R}\left(y_r-\pi_r\right)^2.
\label{eq:brier}
\end{equation}
Equation~\eqref{eq:brier}, the Brier score, evaluates whether predicted probabilities match observed frequencies. For continuous trait or frequency forecasts, the study reports error, interval coverage, and residual structure across environmental strata. A model that identifies a signal in one population but fails in held-out data remains exploratory.

\begin{table}[t]
\caption{Evaluation outcomes and permitted interpretation.}
\label{tab:evaluation}
\centering
\footnotesize
\begin{tabular}{L{0.31\columnwidth} L{0.53\columnwidth}}
\toprule
\textbf{Observed outcome} & \textbf{Permitted conclusion}\\
\midrule
Trajectory observed but alternatives untested & Descriptive evolutionary change in declared data; no selection attribution.\\
Selection model predicts held-out data better than neutral/migration alternatives & Selection-supported conditional on scope and model set.\\
Candidate perturbation changes intermediate and trait with controls/rescue & Mechanistically supported pathway in the tested organismal context.\\
Repeated phenotype but distinct candidate loci across lineages & Phenotypic convergence with molecular non-identity; investigate shared developmental constraints.\\
Model result changes under plausible stage/ancestry/Ne assumptions & Assumption-sensitive result; report sensitivity and do not present a stable mechanism.\\
Historical reconstruction has multiple plausible ancestral paths & Competing historical explanations; retain uncertainty rather than choose a deterministic story.\\
\bottomrule
\end{tabular}
\end{table}

\subsection{Replication and data governance}

Independent replication is not merely a repeat of a software command. Population-genetic replication can involve different populations or time intervals; developmental replication can involve a second laboratory, tissue preparation, or homologous-stage atlas; functional replication can involve an independent construct or rescue. The report should identify which layer was replicated and which layers remain single-study evidence. Negative perturbations, failed rescues, and model failures must be versioned and retained because they constrain causal stories.

The workflow has four reproducibility layers: raw biological materials and collection events; curated genomics, phenotypes, and annotations; locked analysis environments and model specifications; and evidence-labeled reports. Sensitive locality, threatened-population, and controlled genomic data require access restrictions. Collections, organismal interventions, and gene perturbations require permits and welfare review appropriate to the lineage. None is performed by this design-only work.

\begin{table*}[t]
\caption{Failure-oriented decision schedule for an EVOTRACE programme.}
\label{tab:schedule}
\centering
\footnotesize
\begin{tabular}{L{0.16\textwidth} L{0.24\textwidth} L{0.26\textwidth} L{0.22\textwidth}}
\toprule
\textbf{Module} & \textbf{Decision question} & \textbf{Evidence required for advancement} & \textbf{Action when criterion fails}\\
\midrule
Pilot passport and assay & Is the target trait and measurement process stable enough to study? & stage/trait repeatability, sampling coverage, QC, and valid controls & redesign phenotype/stage protocol; do not fit evolutionary mechanism models\\
Population transmission & Does a selected process predict trajectories beyond neutral/demographic alternatives? & calibrated held-out score and replicated trajectory or justified time series & retain descriptive result; collect more temporal or geographic replication\\
Mechanistic bridge & Does a candidate alter a measured developmental/physiological intermediate and trait? & directional evidence, off-target controls, viability check, and orthogonal/rescue result where feasible & retain association only; improve mediator assay or perturbation design\\
Historical synthesis & Are ancestral and comparative routes distinguishable? & explicit alternatives, calibrated phylogeny, homologous trait coding, and comparison lineages & publish competing histories and identify discriminating data\\
Independent evaluation & Does a separate workflow reproduce the inference label? & independent sample, laboratory, or analysis team reproduces the relevant gate result & mark inference provisional and locate disagreement source\\
\bottomrule
\end{tabular}
\end{table*}

Table~\ref{tab:schedule} gives a practical meaning to a failed hypothesis. A neutral model that predicts the data well is not ``nothing happened''; it can rule out an adaptive explanation at the available resolution. A perturbation that changes a developmental intermediate but not the trait can reject a proposed mediation path. A failure to align homologous stages can redirect resources toward the developmental atlas before any comparative gene-expression claim is made.

\section{Discussion}

Natural selection remains essential for explaining adaptive evolution, but a strong account of evolutionary change is plural in mechanism and precise in evidence. Mutation and recombination generate variation; drift and demographic history shape its stochastic fate; gene flow moves variation among populations; ecology determines the context in which fitness is measured; and development maps heritable changes to phenotypes. These components do not compete for a single explanatory slot. They are parts of a causal system whose relative contribution varies by lineage, trait, time scale, and environment.

EVOTRACE provides a disciplined bridge across levels. At the population level, it asks whether a frequency or trait trajectory is better predicted by a selection model than by declared alternatives. At the organismal level, it asks whether the candidate change affects a measured intermediate, trait, and fitness endpoint. At the historical level, it asks whether the proposed route is supported over alternative ancestral paths. The framework avoids a familiar but unproductive dichotomy: either a study discovers one ``gene for'' a complex trait or it claims that mechanism is unknowable. A linked sequence of narrower, testable claims is more informative.

The framework also changes how complex organs are discussed. Eyes, placentas, and other intricate structures need not be treated as exceptions to evolutionary explanation. Their development can involve conserved regulatory modules, co-option, duplication, changed timing, and ecological selection over intermediates. Yet evidence for one of these mechanisms in one component is not evidence for the entire organ's history. Deep homology gives hypotheses about shared generative processes; population data and functional tests determine whether a particular lineage used them in a proposed evolutionary change.

There are practical limits. Effective population size can be difficult to estimate, historical samples may be sparse, selection can be polygenic and context dependent, and developmental assays may be infeasible in nonmodel organisms. Ecological fitness can be expensive to measure and perturbations may have broad pleiotropic effects. EVOTRACE does not remove these limitations. It exposes them in the passport, model alternatives, sensitivity analysis, and claim labels so that a future study can target the limiting evidence rather than overstate its conclusion.

The recommended empirical starting point is a tractable focal trait with replicated populations, an environmental contrast, time-resolved or archival genomic material, and a feasible developmental assay. The project should preregister E0 and E2 model sets, retain both negative and positive results, and add mechanism only after a population-process signal survives alternatives. A later comparative module can test whether the same phenotype reflects parallel regulatory reuse, different molecular routes, or contingency. This modular strategy makes general evolutionary theory grow from calibrated evidence rather than from endpoint stories.

\section{Conclusion}

Organisms evolve because inherited variation is generated and differentially transmitted in populations, under the joint influence of selection, drift, migration, demography, and developmental constraints. Natural selection explains adaptive sorting, but it must be tested against alternative population processes. Complex traits are best explained by evidence that joins population trajectories to developmental intermediates, phenotype, ecology, and fitness. EVOTRACE operationalizes this logic through five stages: scope, variation, transmission, mechanism, and evaluation. It does not claim a new evolutionary result; it supplies a reproducible route by which future studies can make stronger and more transparent claims about how organisms evolve.

\appendices
\section{Minimum Preregistration and Reporting Record}

This fixed record turns the EVOTRACE stages into auditable research commitments. Before collection, a study records the focal lineage, trait ontology, homologous developmental stage, environmental contrast, population boundary, time window, ancestry assumptions, and claim level it intends to evaluate. It also specifies the sampling frame, inclusion and exclusion rules, collection effort, repeated-measure structure, genotype quality thresholds, phenotype calibration procedure, and the treatment of missingness. These fields prevent a measured endpoint from silently changing into a different trait, population, or developmental comparison after results are known.

For transmission, the record lists the primary selection-plus-environment model and the neutral demographic, migration, mutation, and recombination alternatives; it also states the held-out unit, model-comparison metric, and sensitivity ranges for effective population size, ancestry, migration, recombination, and stage alignment. For mechanism, it identifies the candidate intermediate, permitted perturbation or orthogonal evidence, viability and off-target controls, and the result that would downgrade a pathway to association only. For history, it documents the phylogeny, outgroups, trait-coding rule, ancestral-state uncertainty, and the competing routes to be compared.

Every release should preserve raw-material and collection identifiers, data and annotation versions, analysis environment, random seeds where applicable, model specifications, negative results, and deviations from the preregistered record. Reports should state whether each inference stopped at a descriptive, selection-supported, mechanistically supported, or historical label, and why. Sensitive locality, threatened-population, and controlled genomic information may require restricted access, but access limitations and the resulting verification boundary must be declared. This appendix specifies a reporting ledger only; it does not register a completed empirical study or supply observations for one.

\subsection{Deviation and Evidence Ledger}

A deviation is not merely a prose note. The ledger records its date, trigger, affected EVOTRACE stage, changed variable or model, whether the change was made before or after inspection of a relevant outcome, and the anticipated direction of risk. For example, replacing a planned developmental time point with a more convenient specimen is recorded as a stage-alignment change; adding a migration model after a selection signal is observed is recorded as a transmission-model expansion. The final report presents the original and updated specifications, explains the practical reason for the change, and reports whether the inference label changes under both versions. This preserves a usable distinction between adaptation of a protocol to a field constraint and outcome-driven reinterpretation.

Each material claim receives an evidence row with five linked identifiers: the population or lineage record; the genotype, phenotype, or developmental intermediate; the ecological or laboratory condition; the comparison model or control; and the release version containing code and data provenance. The row names the strongest permitted label and its limiting alternative. A selection-supported row therefore points to the neutral, demographic, and migration fits it survived; a mechanism row points to the matched-stage assay, perturbation or orthogonal test, viability evidence, and off-target control; a historical row points to the phylogenetic calibration and alternate ancestral routes. Evidence that cannot be re-linked to these identifiers may remain useful background but cannot promote an evolutionary explanation.

The release decision is made separately for every stage. E1 may be released as a quality-controlled measurement resource even when E2 is inconclusive. E2 may support a conditional population-process account while E3 remains untested. E3 can reject a proposed mediator even if the population trajectory is robust. E4 records replication, calibration, and sensitivity rather than silently converting partial confirmation into generality. A transparent report preserves lower-level observations and identifies the unresolved next experiment. These are future design commitments, not present results.

\balance
\begin{thebibliography}{99}

\bibitem{wiser2013}
M. J. Wiser, N. Ribeck, and R. E. Lenski, ``Long-term dynamics of adaptation in asexual populations,'' \emph{Science}, vol. 342, no. 6164, pp. 1364--1367, 2013, doi: 10.1126/science.1243357.

\bibitem{blount2008}
Z. D. Blount, C. Z. Borland, and R. E. Lenski, ``Historical contingency and the evolution of a key innovation in an experimental population of \emph{Escherichia coli},'' \emph{Proceedings of the National Academy of Sciences}, vol. 105, no. 23, pp. 7899--7906, 2008, doi: 10.1073/pnas.0803151105.

\bibitem{shubin2009}
N. Shubin, C. Tabin, and S. Carroll, ``Deep homology and the origins of evolutionary novelty,'' \emph{Nature}, vol. 457, pp. 818--823, 2009, doi: 10.1038/nature07891.

\bibitem{erwindavidson2009}
D. H. Erwin and E. H. Davidson, ``The evolution of hierarchical gene regulatory networks,'' \emph{Nature Reviews Genetics}, vol. 10, pp. 141--148, 2009, doi: 10.1038/nrg2499.

\bibitem{barrettschluter2008}
R. D. H. Barrett and D. Schluter, ``Adaptation from standing genetic variation,'' \emph{Trends in Ecology \& Evolution}, vol. 23, no. 1, pp. 38--44, 2008, doi: 10.1016/j.tree.2008.08.008.

\bibitem{kaessmann2010}
H. Kaessmann, ``Origins, evolution, and phenotypic impact of new genes,'' \emph{Genome Research}, vol. 20, no. 10, pp. 1313--1326, 2010, doi: 10.1101/gr.101386.109.

\bibitem{nielsen2005}
R. Nielsen, ``Molecular signatures of natural selection,'' \emph{Annual Review of Genetics}, vol. 39, pp. 197--218, 2005, doi: 10.1146/annurev.genet.39.073003.112420.

\bibitem{kimura1983}
M. Kimura, \emph{The Neutral Theory of Molecular Evolution}. Cambridge, U.K.: Cambridge Univ. Press, 1983.

\bibitem{crespsemen2004}
B. J. Crespi and C. A. D. Semeniuk, ``Parent-offspring conflict in the evolution of vertebrate reproductive mode,'' \emph{The American Naturalist}, vol. 163, no. 5, pp. 635--653, 2004, doi: 10.1086/382734.

\bibitem{kawecki2012}
T. J. Kawecki \emph{et al.}, ``Experimental evolution,'' \emph{Trends in Ecology \& Evolution}, vol. 27, no. 10, pp. 547--560, 2012, doi: 10.1016/j.tree.2012.06.001.

\bibitem{courtier2015}
V. Courtier-Orgogozo, ``Replaying the tape of life in the twenty-first century,'' \emph{Interface Focus}, vol. 5, no. 6, 20150057, 2015, doi: 10.1098/rsfs.2015.0057.

\bibitem{stern2008}
D. L. Stern and V. Courtier-Orgogozo, ``The loci of evolution: How predictable is genetic evolution?'' \emph{Evolution}, vol. 62, no. 9, pp. 2155--2177, 2008, doi: 10.1111/j.1558-5646.2008.00450.x.

\bibitem{fave2015}
M.-J. Favé \emph{et al.}, ``Past climate change on Sky Islands drives novelty in a core developmental gene network and its phenotype,'' \emph{BMC Evolutionary Biology}, vol. 15, art. 183, 2015, doi: 10.1186/s12862-015-0448-4.

\bibitem{jablonski2017}
D. Jablonski, ``Approaches to macroevolution: 1. General concepts and origin of variation,'' \emph{Evolutionary Biology}, vol. 44, no. 4, pp. 427--450, 2017, doi: 10.1007/s11692-017-9420-0.

\end{thebibliography}

\end{document}
